\chapter{Fibre}

\section*{Definition and Overview}
Fibre is the part of plant foods that the body cannot digest or absorb. It passes through the digestive system largely intact, and although it provides no energy, it plays essential roles in bowel health, blood sugar control, cholesterol management, and weight regulation. There are two main types: soluble fibre, which dissolves in water to form a gel and helps lower cholesterol and steady blood sugar, and insoluble fibre, which adds bulk to stool and prevents constipation. Most Australians eat only about half the recommended amount of fibre. Increasing fibre through whole foods is one of the simplest and most powerful dietary changes for long-term health. This chapter explains the types of fibre, the foods that contain them, and how to increase intake safely.

\section*{Epidemiology}
Fibre intake is low across most of the Australian population. Adults are recommended to eat 25--30 grams of fibre a day, but average intakes are closer to 20--25 grams, and many people eat less. Low fibre intake contributes to the very high prevalence of constipation, which affects around 15--20\% of the population, and to higher rates of diverticular disease, haemorrhoids, and bowel cancer. Diets rich in fibre are associated with lower rates of heart disease, type 2 diabetes, and bowel cancer, and increasing population fibre intake is recognised as a priority in preventive health policy.

\section*{Aetiology and Pathophysiology}
\begin{itemize}
    \item \textbf{Soluble fibre:} found in oats, barley, legumes, fruits such as apples and citrus, and psyllium; it dissolves in water, forms a gel that slows digestion, lowers LDL cholesterol, and steadies blood sugar
    \item \textbf{Insoluble fibre:} found in wholegrains, wheat bran, nuts, and the skins of vegetables and fruit; it does not dissolve, adds bulk to stool, and speeds the passage of waste
    \item \textbf{Mechanism of action:} fibre increases stool bulk and softness, feeds beneficial gut bacteria, and produces short-chain fatty acids that nourish the bowel lining and reduce inflammation
    \item \textbf{Cholesterol effects:} soluble fibre binds cholesterol in the gut so more is excreted, lowering blood cholesterol
    \item \textbf{Sugar control:} the gel formed by soluble fibre slows the absorption of sugar, reducing blood glucose spikes
    \item \textbf{Satiety:} fibre makes food bulky and slows emptying of the stomach, increasing fullness and helping weight control
    \item \textbf{Gut bacteria:} fibre is the main food source for the healthy bacteria of the gut microbiome, and a fibre-poor diet starves them
\end{itemize}

\section*{Clinical Features}
\begin{itemize}
    \item With adequate fibre: regular, soft, easy bowel movements, usually once a day to three times a week
    \item With inadequate fibre: constipation, hard stools, straining, and infrequent bowel movements
    \item Improved fullness and satisfaction after meals
    \item Better blood sugar control in people with diabetes
    \item Lower cholesterol on blood tests
    \item Symptoms of too-rapid increase: bloating, wind, cramping, and loose stools
    \item No single symptom proves adequacy; bowel regularity and overall dietary pattern are the guides
\end{itemize}

\section*{Differential Diagnosis}
Digestive symptoms sometimes blamed on fibre need sorting out.
\begin{itemize}
    \item Constipation from many causes, including low fibre, low fluids, inactivity, medicines, and medical conditions such as hypothyroidism
    \item Irritable bowel syndrome, in which some people tolerate certain fibres poorly
    \item Coeliac disease and other causes of malabsorption
    \item Bowel cancer and other serious conditions, which can present with changed bowel habit
    \item Diverticular disease and haemorrhoids, which are associated with long-term low fibre intake
    \item The need to increase fibre gradually, since sudden increases cause temporary wind and bloating in everyone
\end{itemize}

\section*{When to Seek Medical Care}
Most people can increase fibre safely without medical advice, but a doctor should be consulted for persistent constipation, blood in the stool, unexplained weight loss, or a change in bowel habit lasting more than a few weeks. People with diabetes adjusting fibre to manage sugar levels should do so with guidance.

\textbf{Call 000 (or local emergency services) for:}
\begin{itemize}
    \item Severe abdominal pain, especially with vomiting and inability to pass wind or stool, which may indicate bowel obstruction
    \item Sudden severe abdominal pain, which may indicate a perforation or other emergency
    \item Significant rectal bleeding or collapse
    \item Any emergency symptoms
\end{itemize}

\section*{Diagnosis and Investigations}
\begin{itemize}
    \item \textbf{Dietary review:} assessment of usual fibre intake against recommendations
    \item \textbf{Bowel symptom history:} frequency, consistency, straining, and any blood
    \item \textbf{Blood tests:} for iron deficiency, thyroid disease, or coeliac disease where indicated
    \item \textbf{Investigation of changed bowel habit:} colonoscopy or other tests for people with concerning symptoms or risk factors
    \item \textbf{No routine testing:} for most people, increasing fibre is a lifestyle change rather than a medical treatment
\end{itemize}

\section*{Management}
\begin{itemize}
    \item \textbf{Choose wholegrains:} replace white bread, white rice, and refined cereals with wholegrain versions, oats, and brown rice
    \item \textbf{Eat more legumes:} lentils, chickpeas, beans, and peas are among the richest fibre sources
    \item \textbf{Fruit and vegetables:} aim for at least five serves a day, with skins left on where possible
    \item \textbf{Nuts and seeds:} add almonds, walnuts, chia seeds, and flaxseed to meals and snacks
    \item \textbf{Increase gradually:} add fibre slowly over a few weeks, with plenty of fluid, to avoid wind and bloating
    \item \textbf{Fibre supplements:} products such as psyllium can help when dietary intake cannot be increased
    \item \textbf{Match fluid intake:} fibre works best with adequate water; increasing fibre without fluids can worsen constipation
    \item \textbf{Specific conditions:} people with irritable bowel syndrome may tolerate soluble fibre better than insoluble fibre and should seek tailored advice
\end{itemize}

\section*{Complications}
\begin{itemize}
    \item Constipation and haemorrhoids from chronic low fibre intake
    \item Diverticular disease, which is more common in low-fibre diets
    \item Increased risk of bowel cancer, heart disease, and type 2 diabetes with lifelong low fibre intake
    \item Wind, bloating, and cramps if fibre is increased too quickly
    \item Potential worsening of symptoms in irritable bowel syndrome without individualised advice
\end{itemize}

\section*{Prognosis and Outcomes}
Increasing fibre produces measurable benefits within weeks: improved bowel regularity, better fullness, and, over months, lower cholesterol and steadier blood sugar. Long-term high fibre intake is associated with substantially lower rates of heart disease, stroke, diabetes, and bowel cancer, and with a healthier gut microbiome. The benefits are greatest when fibre comes from whole foods rather than supplements, and when the change is sustained. For most people, reaching the recommended 25--30 grams a day is achievable with modest, gradual dietary changes and brings lasting health gains.

\section*{Prevention and Self-Care}
\begin{itemize}
    \item Build meals around wholegrains, legumes, vegetables, and fruit
    \item Aim for 25--30 grams of fibre daily, increasing gradually from current levels
    \item Drink plenty of water alongside higher fibre intake
    \item Read food labels and choose products with higher fibre per serve
    \item Keep the skins on fruit and vegetables where edible
    \item Use breakfast as an easy opportunity: wholegrain cereal or oats with fruit
    \item If using a fibre supplement, introduce it gradually and follow instructions
\end{itemize}

\section*{Sources and Further Reading}
The following reputable sources provide further information for healthcare students and patients seeking reliable, current advice about dietary fibre.
\begin{itemize}
    \item National Health and Medical Research Council. \textit{Australian Dietary Guidelines.} https://www.eatforhealth.gov.au/
    \item Healthdirect Australia. \textit{Dietary fibre.} https://www.healthdirect.gov.au/fibre
    \item Dietitians Australia. \textit{Fibre and your health.} https://dietitiansaustralia.org.au/
    \item Better Health Channel. \textit{Fibre in food.} https://www.betterhealth.vic.gov.au/
    \item NHS. \textit{How to get more fibre in your diet.} https://www.nhs.uk/live-well/eat-well/
    \item World Health Organization. \textit{Carbohydrate and fibre guidance.} https://www.who.int/
\end{itemize}
